% !TEX program = pdflatex
% Exercise 5 adapts REP-02 from the instructor's midterm exercise bank.
% Exercise 6 uses original prompts in the oral style of the local Finance Games bank.
\documentclass[10pt,letterpaper]{article}
\usepackage[T1]{fontenc}
\usepackage{lmodern,microtype}
\usepackage[margin=0.7in,includeheadfoot,headheight=13pt,headsep=10pt,footskip=22pt]{geometry}
\usepackage{amsmath,amssymb,mathtools,enumitem,xcolor,fancyhdr,lastpage,hyperref}
\definecolor{ink}{HTML}{193B56}
\definecolor{muted}{HTML}{52616C}
\hypersetup{hidelinks,pdftitle={Rapid Fire - 60 Short Questions - Lectures 1-11},
 pdfauthor={Nathan Sauldubois},pdfsubject={Stochastic calculus and option pricing: self-check questions}}
\pagestyle{fancy}\fancyhf{}
\fancyhead[L]{\small\color{ink}\textbf{STOCHASTIC CALCULUS \& OPTION PRICING}}
\fancyhead[R]{\small\color{muted}Rapid fire | Lectures 1--11}
\fancyfoot[L]{\small\color{muted}NYU Tandon | Fall 2026}
\fancyfoot[R]{\small\color{muted}\thepage\,/\,\pageref{LastPage}}
\renewcommand{\headrulewidth}{0.35pt}
\setlength{\parindent}{0pt}\setlength{\parskip}{3pt}
\setlength{\abovedisplayskip}{5pt}\setlength{\belowdisplayskip}{5pt}
\setlength{\abovedisplayshortskip}{3pt}\setlength{\belowdisplayshortskip}{4pt}
\newcommand{\dd}{\,\mathrm{d}}
\newcommand{\E}{\mathbb{E}}
\newcommand{\Pbb}{\mathbb{P}}
\newcommand{\Q}{\mathbb{Q}}
\newcommand{\F}{\mathcal{F}}
\newcommand{\round}[2]{\par\vspace{7pt}{\large\bfseries\color{ink}Exercise #1\quad #2\par}\vspace{2pt}}
\newenvironment{drill}[1]{\begin{enumerate}[label=\textbf{#1.\arabic*},
 leftmargin=2.7em,labelsep=0.6em,itemsep=4pt,topsep=4pt,parsep=0pt]}{\end{enumerate}}
\raggedbottom
\begin{document}
\fontsize{10.7}{13}\selectfont
{\LARGE\bfseries\color{ink}Rapid fire\par}
{\color{muted}60 short questions to check what you can do without your notes.\par}
Most answers need one line or a few equations. Exercise 5 is a guided replication argument.
Mark any item you cannot justify, then return to it after reviewing the course.
Time is in years. Unless stated otherwise, rates are constant, stocks pay no dividends and trading is frictionless.

\round{1}{Quick applications of It\^o: one and two states}
Give the \emph{full differential}, including the drift and every Brownian term.
Here $W$ is a standard Brownian motion. For Questions 1.6--1.8,
$\dd\langle W^1,W^2\rangle_t=\tfrac12\dd t$.
\begin{drill}{1}
\item Compute $\dd(W_t^2)$.
\item Compute $\dd(W_t^3)$.
\item Compute $\dd(e^{2W_t})$.
\item If $\dd X_t=\mu X_t\dd t+\sigma X_t\dd W_t$, with $X_0>0$, compute $\dd\log X_t$.
\item If $\dd X_t=\dd t+2\dd W_t$, compute $\dd(tX_t^2)$.
\item Compute $\dd(W_t^1W_t^2)$. Which term disappears when the two drivers are independent?
\item Let $\dd X_t=\dd t+2\dd W_t^1$ and $\dd Y_t=-2\dd t+3\dd W_t^2$.
Compute $\dd(X_tY_t)$.
\item With the same $X,Y$, compute $\dd[(X_t+Y_t)^2]$.
\end{drill}

\round{2}{Girsanov: change the measure, recalculate the drift}
On $[0,T]$, $W$ is Brownian under $\Pbb$, and $\lambda$ is bounded and deterministic. Define
\[
 \frac{\dd\Q}{\dd\Pbb}\bigg|_{\F_T}
 =Z_T=\exp\!\left(-\int_0^T\lambda_u\dd W_u
                  -\tfrac12\int_0^T\lambda_u^2\dd u\right).
\]
\begin{drill}{2}
\item Give a sufficient condition ensuring that $Z$ is a true density martingale.
Does the stated $\lambda$ satisfy it?
\item Write the $\Q$-Brownian motion $W^{\Q}$ and express $\dd W$ in terms of $\dd W^{\Q}$.
\item Under $\Pbb$, $\dd X_t=3\dd t+2\dd W_t$. For $\lambda_t=1$,
give its drift and diffusion under $\Q$.
\item For the same $X$ with deterministic $X_0=x_0$, now take $\lambda_t=t$.
Give the $\Q$ dynamics and the distribution of $X_T$.
\item A stock has $\dd S/S=0.08\dd t+0.20\dd W$ under $\Pbb$, while $r=0.03$.
Choose $\lambda$ for risk-neutral pricing; give the $\Q$ drift of $\log S$.
\item Finally, use the density built from $W^1$ with $\lambda=0.4$.
Under $\Pbb$, $\dd Y/Y=0.06\dd t+0.30\dd W^2$, with correlation $\rho=1/2$.
Find $W^{2,\Q}$ and the relative drift of $Y$ under $\Q$. Does the correlation change?
\end{drill}

\newpage
\round{3}{Martingale or not?}
Use the augmented natural Brownian filtration on a finite interval $[0,T]$ under $\Pbb$.
For each process, answer \textbf{yes or no} and justify it. Check adaptedness and integrability;
use a conditional expectation or It\^o's formula. Here $\theta\ne0$ is constant.
\begin{drill}{3}
\item $M_t=W_t$.
\item $M_t=W_t^2$.
\item $M_t=W_t^2-t$.
\item $M_t=W_t^3-3tW_t$.
\item $M_t=\exp(\theta W_t-\tfrac12\theta^2t)$.
\item $M_t=\exp(\theta W_t)$.
\item $M_t=\int_0^t W_u\dd W_u$.
\item $M_t=\int_0^t W_u^2\dd u$.
\item $M_t=W_T$ for every $0\leq t\leq T$.
\item $M_t=W_{t/2}$. Is being adapted with constant expectation sufficient?
\end{drill}

\round{4}{Quote checks: is there an arbitrage?}
Each question is an \textbf{independent market}. All options are European on the same stock,
with $S_0=100$, $T=1$, $r=0$ and no dividends; borrowing, lending and short selling are allowed.
Write $C(K)$ and $P(K)$ for option prices at strike $K$.
Ignore bid--ask spreads here.

For each set, identify a violated no-arbitrage condition, or state that the quoted checks pass.
If you find an arbitrage, specify the positions, the initial receipt and a nonnegative terminal payoff.
\begin{drill}{4}
\item $C(90)=8$.
\item $C(100)=9$, $P(100)=12$.
\item $C(90)=14$, $C(110)=16$.
\item $C(90)=24$, $C(110)=2$.
\item $C(80)=25$, $C(100)=20$, $C(130)=5$. Use the actual strike spacings.
\item $C(90)=15$, $C(100)=9$, $C(110)=4$.
Check the bounds, monotonicity, slope bounds and convexity, including the point $C(0)=S_0$.
\end{drill}

\newpage
\round{5}{Replicate using the instrument that is actually traded}
Only a bank account $B$ and an asset $X$ can be traded. Under $\Pbb$,
\[
 B_t=e^{rt},\qquad
 \dd X_t=\mu X_t\dd t+\sigma X_t\dd W_t,
 \qquad X_0>0,\quad\sigma>0.
\]
The constants are $r,\mu,\sigma$. The observed variable is $S_t=\sqrt{X_t}$;
there is no separately traded asset with price $S_t$.
The contract pays \textbf{the cash amount $S_T$} at $T$ and makes no earlier payments.
Write its price as $v(t,S_t)$, and let $\phi_t,\beta_t$ be units of $X,B$.
Assume no arbitrage, sufficient smoothness, and the integrability and admissibility needed for replication.

Derive the answer from the trading equations. Keep a number of units distinct from a cash amount.
\begin{drill}{5}
\item Name the traded instruments. Can $S$ itself appear as an independently chosen holding?
\item Apply It\^o to $S=\sqrt X$ and give its dynamics under $\Pbb$.
\item Write the portfolio value $V$ and its self-financing gain equation in terms of $X,B$.
\item Apply It\^o to $v(t,S_t)$. Give the drift and Brownian coefficient.
\item Match Brownian terms to find $\phi$ in terms of $v_s$ and $S$.
Why is holding $v_s$ units of $X$ generally incorrect?
\item Express the cash amount $\beta B$ in terms of $v,S,v_s$.
\item Match drifts to derive the pricing PDE and its terminal condition.
Where does the historical drift $\mu$ disappear?
\item Try $v(t,s)=A(t)s$. Derive the ODE and terminal value for $A$.
\item Solve the ODE and give the price $v(t,s)$ explicitly.
\item Give the resulting $\phi$ and $\beta B$. Check their self-financing gains and
the value at maturity.
\end{drill}

\newpage
\round{6}{Interview round: explain it in one sentence}
Answer directly, then defend the answer with a formula, a sign or a counterexample.
Use the standard BS benchmark unless the question changes the model.
\begin{drill}{6}
\item A European call has spot 100, strike 100 and one year to expiry.
What inputs are still missing for a BS price? Do you need the historical drift $\mu$?
\item With $S=K=100$, $r=0$ and $\sigma=20\%$, estimate the price of a call expiring
in one business day, taken as $1/252$ year.
\item You are long one call. How do you make the position delta neutral with stock?
If spot rises with the other inputs fixed, do you buy or sell more stock to rebalance?
\item As expiry approaches at the money, what happens to gamma and vega?
\item A call and a put have the same strike and maturity. Which has the larger BS vega, and why?
\item For the position ``long call, short put'' at the same strike and maturity,
give the terminal payoff, gamma and vega.
\item Does higher BS volatility raise both vanilla call and put prices?
Must it also raise the price of a nonnegative call-spread payoff?
\item Holding the other BS inputs fixed, how does a higher interest rate affect a European call and put?
\item Which is worth more: an American or European call? Give the standard condition for equality.
\item Can you price an American option by selecting the best discounted exercise payoff
on each simulated path after observing the whole path?
\item You buy a call at 20\% implied volatility and delta hedge daily.
Realized volatility is 30\%. Is a positive terminal profit guaranteed?
\item Equity implied volatility is higher at lower strikes. Does this permit a call and put
at the \emph{same} strike and expiry to have different parity-consistent implied volatilities?
\item Distinguish implied volatility from the local-volatility coefficient $\sigma_{\rm loc}(t,s)$.
\item A stochastic-volatility model adds an independent variance shock.
Can stock and cash generally replicate every claim?
\item Daily returns have a 1\% standard deviation and are independent with equal variance.
What is the annualized volatility over 252 days? What if every daily return is deterministically $+1\%$?
\item At fixed log-forward moneyness $k=\log(K/F_{0,T})$, where $F_{0,T}=S_0e^{rT}$,
3-month implied volatility is 30\% and 1-year implied volatility is 20\%.
Does this pair violate the total-variance calendar check?
\item The same payoff is quoted bid/ask at 8.8/9.2 on one venue and 9.0/9.4 on another.
Do the different mid-prices give an executable arbitrage?
\item You multiply the number of independent MC paths by four. What happens to the standard error?
Does this reduce Euler time-discretization bias?
\item A simulation uses \texttt{h*Z} for a Brownian increment, with $Z\sim N(0,1)$.
What is wrong, and what is the correct scaling?
\item You already own a stock. Does selling a put protect you against a stock-price fall?
What put position creates a protective floor instead?
\end{drill}
\end{document}
